\section{Постановка задачи}
\subsection{Данные, предположения и ограничения}
Даны запросы и эталоны $(x,y)\sim P$, два фиксированных классификатора
$M_w,M_s$ и выборка парных исходов их работы.
Обозначим $W=\mathbf1\{M_w(x)\text{ правилен}\}$,
$S=\mathbf1\{M_s(x)\text{ правилен}\}$ и $D=S-W\in\{-1,0,1\}$.
Здесь $D=1$ --- спасение, $D=-1$ --- вредная замена,
$D=0$ --- одинаковая правильность ответов. Предположение о независимости
ошибок моделей не требуется; большая accuracy второй модели не постулируется.

Сначала вызывается $M_w$. Роутеру доступны признаки
$z=\phi(x,M_w(x))$, которые можно получить до второго вызова.
Текущий ответ $M_s$ и эталон при решении неизвестны. При эскалации
второй ответ заменяет первый без отдельного судьи или объединения ответов.
Правильность обеих моделей используется только как обучающая разметка
или для оценки качества.

Эксперты обучаются на отдельной части данных и затем фиксируются.
Router-fit, policy-dev и test не пересекаются по объектам с этой частью
или друг с другом. Все признаки преобразуются по параметрам, оценённым
только на router-fit. В базовом опыте train/dev/test роутера представляют
естественный поток. При целевом отборе router-train его распределение
может отличаться от $P$; поправка на такой отбор должна учитываться
отдельно. Новые запросы рассматриваются как выборка из целевого потока,
а возможная семантическая зависимость запросов ограничивает перенос
обычных IID-оценок неопределённости.

\subsection{Полезность, стоимость и внешние критерии}
Условная полезность эскалации:
\begin{equation}\label{eq:gain}
\begin{aligned}
g(z)&=\mathbb E[D\mid z]\\
&=P(W=0,S=1\mid z)-P(W=1,S=0\mid z)\\
&=P(S=1\mid z)-P(W=1\mid z).
\end{aligned}
\end{equation}
Для решения $r(z)\in\{0,1\}$ внешний критерий --- accuracy каскада:
\begin{equation}\label{eq:quality}
 Q(r)=\mathbb E[W+r(z)D],\qquad
 \mathbb E[c_w+c_{\phi,r}+r(z)c_s]\le B.
\end{equation}
Здесь $c_w$ --- первый вызов, $c_{\phi,r}$ --- признаки и роутер,
$c_s$ --- дополнительный второй вызов. В базовом опыте вместо
измеренного денежного бюджета используется одинаковая для всех методов
квота: ровно $k=\lfloor bN\rfloor$ вызовов на пакет из $N$ запросов.
Выбираются top - $k$ оценок с общим разрешением равных scores.
Это пакетный отбор, а не проверенный порог для online-потока.

Профиль принятого решения относительно противоположного действия
описывается $Z_r=(2r-1)D$ и таблицей
\begin{equation}
 T_{v,a}=\sum_{i=1}^{N}\mathbf1\{Z_{r,i}=v,r_i=a\},
 \qquad v\in\{-1,0,1\},\quad a\in\{0,1\}.
\end{equation}
При эскалации строка $+1$ означает спасение, $-1$ --- вред;
без эскалации $+1$ означает предотвращённый вред, $-1$ --- упущенное
спасение. Нейтральный исход объединяет два правильных и два ошибочных
ответа. При числе спасений $R$ и вредных замен $H$ эмпирический прирост
accuracy равен $(R-H)/N$. Поэтому таблица раскрывает компоненты
агрегированной метрики, но не заменяет проверку их различий.

\subsection{Внутренний критерий и процедура выбора}
Для каждого способа обучения параметры оптимизируются на router-fit,
сила регуляризации выбирается по MSE наблюдаемого $D$ на policy-dev.
Критерий настройки отделён от внешней accuracy при фиксированном $b$.
Для вероятностных голов используется log-loss; отсутствующий класс
предсказывается с нулевой вероятностью, единственный --- константой.
В прямом варианте ищется линейная функция $f_\theta(z)$; intercept
не штрафуется. Внутренняя задача:
\begin{equation}\label{eq:argmin}
\widehat\theta_\lambda\in\arg\min_\theta
\left\lbrace
\frac1n\sum_{i=1}^{n}(D_i-f_\theta(z_i))^2+
\lambda\|\theta_{\mathrm{coef}}\|_2^2
\right\rbrace.
\end{equation}

Основная исследовательская задача --- выбор способа обучения
$a\in\mathcal A$ по разрешённым характеристикам данных, пары и бюджета.
Для конфигурации $\tau$ обозначим обученный роутер $r_{\tau,a}$,
его качество $Q_\tau(r_{\tau,a})$ и descriptors $d_\tau$.
Рекомендация $\pi(d_\tau)$ не использует test; на новой задаче допускаются
обучение и настройка роутера на её разрешённых train/dev. Контроль
переноса требует целиком отложенных задач или пар. При
$a_\tau^*\in\arg\max_{a\in\mathcal A}Q_\tau(r_{\tau,a})$
задача выбора записывается как
\begin{equation}\label{eq:selection}
\pi^* \in \mathop{\mathrm{arg\,min}}\limits_{\pi\in\Pi}
\mathbb{E}_\tau
\bigl[\,Q_\tau(r_{\tau,a_\tau^*})-
Q_\tau(r_{\tau,\pi(d_\tau)})\,\bigr]
\end{equation}
при сопоставимом бюджете парной разметки, включая настройку.
